We introduced \textbf{RetailBench}, a data-grounded benchmark for evaluating long-horizon, tool-using LLM agents in single-store retail operation. RetailBench exposes agents to a closed-loop decision process involving pricing, replenishment, supplier selection, shelf assortment, customer feedback, external events, and financial constraints.

Our results show that current LLM agents still struggle to sustain effective long-horizon operation. Even when some models survive the full evaluation horizon, they remain well below the heuristic reference in net worth and sales outcomes. Failure analysis identifies three recurring bottlenecks: incomplete evidence acquisition, surface-level decision making, and lack of a consistent long-horizon policy. 

These findings highlight that reliable deployment in realistic, complex production environments requires agents to go beyond isolated task solving. Future agents must select relevant evidence, reason over it carefully, translate it into effective actions, maintain coherent stateful policies, and adapt to feedback over extended horizons. Failure in any of these capabilities can degrade agent performance in long-horizon environments. This provides a concrete direction for improving both agent architectures and the underlying LLMs.